% ============================================================
% Question Bank: ch01-03 Finding the Constant of Proportionality
% Topic Title: Finding the Constant of Proportionality
% CCSS: 7.RP.A.2.B
% Questions: 30 (18 MC + 7 SA + 5 Graphical)
% ============================================================

% --- Multiple Choice Questions (18) ---

\begin{practiceQuestion}{1-3-q01}{mc}
\begin{questionText}
What is the constant of proportionality ($k$) for the table below?

\begin{center}
\begin{tabular}{|c|c|}
\hline $x$ & $y$ \\ \hline $3$ & $12$ \\ $5$ & $20$ \\ $8$ & $32$ \\ \hline
\end{tabular}
\end{center}
\end{questionText}
\choiceA{$3$}
\choiceB{$4$}
\choiceC{$8$}
\choiceD{$12$}
\correctAnswer{B}
\explanation{$k = \frac{y}{x} = \frac{12}{3} = 4$. Check: $\frac{20}{5} = 4$ and $\frac{32}{8} = 4$.}
\end{practiceQuestion}

\begin{practiceQuestion}{1-3-q02}{mc}
\begin{questionText}
In a proportional relationship, the constant of proportionality $k$ equals:
\end{questionText}
\choiceA{$x \div y$}
\choiceB{$y \div x$}
\choiceC{$x \times y$}
\choiceD{$y - x$}
\correctAnswer{B}
\explanation{The constant of proportionality is $k = \frac{y}{x}$, which is the unit rate.}
\end{practiceQuestion}

\begin{practiceQuestion}{1-3-q03}{mc}
\begin{questionText}
A proportional graph passes through the point $(4, 18)$. What is $k$?
\end{questionText}
\choiceA{$4$}
\choiceB{$4.5$}
\choiceC{$14$}
\choiceD{$18$}
\correctAnswer{B}
\explanation{$k = \frac{y}{x} = \frac{18}{4} = 4.5$.}
\end{practiceQuestion}

\begin{practiceQuestion}{1-3-q04}{mc}
\begin{questionText}
Bananas cost a fixed price per pound. The table shows the relationship. What is $k$?

\begin{center}
\begin{tabular}{|c|c|}
\hline \textbf{Pounds} & \textbf{Cost (\$)} \\ \hline $2$ & $\$1.30$ \\ $5$ & $\$3.25$ \\ $8$ & $\$5.20$ \\ \hline
\end{tabular}
\end{center}
\end{questionText}
\choiceA{$\$0.55$ per pound}
\choiceB{$\$0.65$ per pound}
\choiceC{$\$1.30$ per pound}
\choiceD{$\$1.54$ per pound}
\correctAnswer{B}
\explanation{$k = \frac{1.30}{2} = 0.65$. Check: $\frac{3.25}{5} = 0.65$ and $\frac{5.20}{8} = 0.65$. Bananas cost $\$0.65$ per pound.}
\end{practiceQuestion}

\begin{practiceQuestion}{1-3-q05}{mc}
\begin{questionText}
A car travels at a constant speed. It goes $150$ miles in $2.5$ hours. What is the constant of proportionality?
\end{questionText}
\choiceA{$30$ mph}
\choiceB{$60$ mph}
\choiceC{$75$ mph}
\choiceD{$150$ mph}
\correctAnswer{B}
\explanation{$k = \frac{150}{2.5} = 60$ miles per hour.}
\end{practiceQuestion}

\begin{practiceQuestion}{1-3-q06}{mc}
\begin{questionText}
On a proportional graph, the point $(1, r)$ represents:
\end{questionText}
\choiceA{The $x$-intercept}
\choiceB{The constant of proportionality}
\choiceC{The slope of the line}
\choiceD{Both B and C}
\correctAnswer{D}
\explanation{At $x = 1$, $y = k \times 1 = k$. So the point $(1, r)$ gives the constant of proportionality, which is also the slope.}
\end{practiceQuestion}

\begin{practiceQuestion}{1-3-q07}{mc}
\begin{questionText}
The table shows hours worked and money earned. What is $k$ and what does it mean?

\begin{center}
\begin{tabular}{|c|c|}
\hline \textbf{Hours} & \textbf{Earnings (\$)} \\ \hline $\frac{1}{2}$ & $\$6$ \\ $1$ & $\$12$ \\ $3$ & $\$36$ \\ \hline
\end{tabular}
\end{center}
\end{questionText}
\choiceA{$k = 6$; earns $\$6$ per half hour}
\choiceB{$k = 12$; earns $\$12$ per hour}
\choiceC{$k = 36$; earns $\$36$ per $3$ hours}
\choiceD{$k = 3$; works $3$ hours per $\$36$}
\correctAnswer{B}
\explanation{$k = \frac{6}{0.5} = 12$ or $\frac{12}{1} = 12$ or $\frac{36}{3} = 12$. The unit rate is $\$12$ per hour.}
\end{practiceQuestion}

\begin{practiceQuestion}{1-3-q08}{mc}
\begin{questionText}
A proportional relationship has $k = 3.5$. If $x = 8$, what is $y$?
\end{questionText}
\choiceA{$11.5$}
\choiceB{$24$}
\choiceC{$28$}
\choiceD{$4.5$}
\correctAnswer{C}
\explanation{$y = kx = 3.5 \times 8 = 28$.}
\end{practiceQuestion}

\begin{practiceQuestion}{1-3-q09}{mc}
\begin{questionText}
Find $k$ from the table.

\begin{center}
\begin{tabular}{|c|c|}
\hline $x$ & $y$ \\ \hline $6$ & $9$ \\ $10$ & $15$ \\ $14$ & $21$ \\ \hline
\end{tabular}
\end{center}
\end{questionText}
\choiceA{$1.5$}
\choiceB{$3$}
\choiceC{$\frac{2}{3}$}
\choiceD{$6$}
\correctAnswer{A}
\explanation{$k = \frac{9}{6} = 1.5$. Check: $\frac{15}{10} = 1.5$ and $\frac{21}{14} = 1.5$.}
\end{practiceQuestion}

\begin{practiceQuestion}{1-3-q10}{mc}
\begin{questionText}
A printer prints $80$ pages in $5$ minutes at a constant rate. What is the constant of proportionality?
\end{questionText}
\choiceA{$5$ pages/min}
\choiceB{$16$ pages/min}
\choiceC{$75$ pages/min}
\choiceD{$80$ pages/min}
\correctAnswer{B}
\explanation{$k = \frac{80}{5} = 16$ pages per minute.}
\end{practiceQuestion}

\begin{practiceQuestion}{1-3-q11}{mc}
\begin{questionText}
The graph of a proportional relationship passes through the points $(0, 0)$ and $(5, 12)$. What is $k$?
\end{questionText}
\choiceA{$2.4$}
\choiceB{$5$}
\choiceC{$7$}
\choiceD{$12$}
\correctAnswer{A}
\explanation{$k = \frac{y}{x} = \frac{12}{5} = 2.4$.}
\end{practiceQuestion}

\begin{practiceQuestion}{1-3-q12}{mc}
\begin{questionText}
A proportional table has $k = \frac{3}{4}$. Which pair could be in the table?
\end{questionText}
\choiceA{$(8, 10)$}
\choiceB{$(8, 6)$}
\choiceC{$(12, 8)$}
\choiceD{$(4, 2)$}
\correctAnswer{B}
\explanation{If $k = \frac{3}{4}$, then $y = \frac{3}{4}x$. For $x = 8$: $y = \frac{3}{4} \times 8 = 6$. So $(8, 6)$ works.}
\end{practiceQuestion}

\begin{practiceQuestion}{1-3-q13}{mc}
\begin{questionText}
The table shows a proportional relationship. Find the missing value.

\begin{center}
\begin{tabular}{|c|c|}
\hline $x$ & $y$ \\ \hline $4$ & $14$ \\ $6$ & $?$ \\ $10$ & $35$ \\ \hline
\end{tabular}
\end{center}
\end{questionText}
\choiceA{$18$}
\choiceB{$20$}
\choiceC{$21$}
\choiceD{$24$}
\correctAnswer{C}
\explanation{$k = \frac{14}{4} = 3.5$. Missing value: $y = 3.5 \times 6 = 21$. Check: $\frac{35}{10} = 3.5$ \checkmark.}
\end{practiceQuestion}

\begin{practiceQuestion}{1-3-q14}{mc}
\begin{questionText}
A recipe uses cups of flour proportional to servings. If $12$ servings need $9$ cups, what is $k$ in cups per serving?
\end{questionText}
\choiceA{$\frac{3}{4}$}
\choiceB{$\frac{4}{3}$}
\choiceC{$3$}
\choiceD{$9$}
\correctAnswer{A}
\explanation{$k = \frac{9}{12} = \frac{3}{4}$ cups per serving.}
\end{practiceQuestion}

\begin{practiceQuestion}{1-3-q15}{mc}
\begin{questionText}
Two students find $k$ from the point $(8, 20)$. Student A says $k = 2.5$. Student B says $k = 0.4$. Who is correct?
\end{questionText}
\choiceA{Student A, because $k = \frac{y}{x}$}
\choiceB{Student B, because $k = \frac{x}{y}$}
\choiceC{Both are correct, depending on which variable is independent}
\choiceD{Neither is correct}
\correctAnswer{A}
\explanation{By convention, $k = \frac{y}{x} = \frac{20}{8} = 2.5$. Student B computed $\frac{x}{y}$, which is the reciprocal — not $k$.}
\end{practiceQuestion}

\begin{practiceQuestion}{1-3-q16}{mc}
\begin{questionText}
A map uses a scale where $1$ cm represents $25$ km. What is the constant of proportionality in km per cm?
\end{questionText}
\choiceA{$0.04$}
\choiceB{$1$}
\choiceC{$25$}
\choiceD{$250$}
\correctAnswer{C}
\explanation{$k = \frac{25 \text{ km}}{1 \text{ cm}} = 25$ km per cm.}
\end{practiceQuestion}

\begin{practiceQuestion}{1-3-q17}{mc}
\begin{questionText}
The table shows a proportional relationship. Find $k$.

\begin{center}
\begin{tabular}{|c|c|}
\hline $x$ & $y$ \\ \hline $\frac{1}{2}$ & $4$ \\ $1$ & $8$ \\ $\frac{3}{2}$ & $12$ \\ \hline
\end{tabular}
\end{center}
\end{questionText}
\choiceA{$4$}
\choiceB{$6$}
\choiceC{$8$}
\choiceD{$12$}
\correctAnswer{C}
\explanation{$k = \frac{4}{\frac{1}{2}} = 4 \times 2 = 8$. Check: $\frac{8}{1} = 8$ and $\frac{12}{\frac{3}{2}} = 12 \times \frac{2}{3} = 8$.}
\end{practiceQuestion}

\begin{practiceQuestion}{1-3-q18}{mc}
\begin{questionText}
Which table does NOT have a constant of proportionality?
\end{questionText}
\choiceA{\begin{tabular}{|c|c|} \hline $x$ & $y$ \\ \hline $2$ & $10$ \\ $4$ & $20$ \\ \hline \end{tabular}}
\choiceB{\begin{tabular}{|c|c|} \hline $x$ & $y$ \\ \hline $3$ & $12$ \\ $6$ & $24$ \\ \hline \end{tabular}}
\choiceC{\begin{tabular}{|c|c|} \hline $x$ & $y$ \\ \hline $4$ & $10$ \\ $8$ & $22$ \\ \hline \end{tabular}}
\choiceD{\begin{tabular}{|c|c|} \hline $x$ & $y$ \\ \hline $5$ & $15$ \\ $10$ & $30$ \\ \hline \end{tabular}}
\correctAnswer{C}
\explanation{A: $k = 5$. B: $k = 4$. C: $\frac{10}{4} = 2.5$ but $\frac{22}{8} = 2.75$ — ratios differ. D: $k = 3$. Only C is not proportional.}
\end{practiceQuestion}

% --- Short Answer Questions (7) ---

\begin{practiceQuestion}{1-3-q19}{sa}
\begin{questionText}
A proportional graph passes through $(6, 15)$. What is $k$?
\end{questionText}
\correctAnswer{$k = 2.5$}
\explanation{$k = \frac{y}{x} = \frac{15}{6} = 2.5$.}
\end{practiceQuestion}

\begin{practiceQuestion}{1-3-q20}{sa}
\begin{questionText}
A store sells ribbon at a constant price per yard. If $3$ yards cost $\$4.50$, what is the constant of proportionality?
\end{questionText}
\correctAnswer{$k = \$1.50$ per yard}
\explanation{$k = \frac{4.50}{3} = 1.50$ dollars per yard.}
\end{practiceQuestion}

\begin{practiceQuestion}{1-3-q21}{sa}
\begin{questionText}
Find the missing value in the proportional table.

\begin{center}
\begin{tabular}{|c|c|}
\hline $x$ & $y$ \\ \hline $5$ & $20$ \\ $9$ & $?$ \\ \hline
\end{tabular}
\end{center}
\end{questionText}
\correctAnswer{$y = 36$}
\explanation{$k = \frac{20}{5} = 4$. So $y = 4 \times 9 = 36$.}
\end{practiceQuestion}

\begin{practiceQuestion}{1-3-q22}{sa}
\begin{questionText}
A runner covers $\frac{3}{4}$ mile in $\frac{1}{4}$ hour at a constant rate. What is $k$ in miles per hour?
\end{questionText}
\correctAnswer{$k = 3$ mph}
\explanation{$k = \frac{\frac{3}{4}}{\frac{1}{4}} = \frac{3}{4} \times 4 = 3$ mph.}
\end{practiceQuestion}

\begin{practiceQuestion}{1-3-q23}{sa}
\begin{questionText}
Two quantities are proportional with $k = 7$. If $y = 56$, what is $x$?
\end{questionText}
\correctAnswer{$x = 8$}
\explanation{$y = kx$ so $56 = 7x$, giving $x = 56 \div 7 = 8$.}
\end{practiceQuestion}

\begin{practiceQuestion}{1-3-q24}{sa}
\begin{questionText}
The table shows a proportional relationship. Find $k$ and the missing value.

\begin{center}
\begin{tabular}{|c|c|}
\hline $x$ & $y$ \\ \hline $3$ & $7.5$ \\ $?$ & $15$ \\ $10$ & $25$ \\ \hline
\end{tabular}
\end{center}
\end{questionText}
\correctAnswer{$k = 2.5$; missing $x = 6$}
\explanation{$k = \frac{7.5}{3} = 2.5$. For the missing row: $15 = 2.5x$ so $x = 15 \div 2.5 = 6$.}
\end{practiceQuestion}

\begin{practiceQuestion}{1-3-q25}{sa}
\begin{questionText}
A factory produces widgets at a constant rate. In $2\frac{1}{2}$ hours it makes $175$ widgets. What is $k$ in widgets per hour?
\end{questionText}
\correctAnswer{$k = 70$ widgets per hour}
\explanation{$k = \frac{175}{2.5} = 70$ widgets per hour.}
\end{practiceQuestion}

% --- Graphical Questions (3 GMC + 2 GSA) ---

\begin{practiceQuestion}{1-3-q26}{gmc}
\begin{questionText}
The graph below shows a proportional relationship. What is the constant of proportionality?

\begin{center}
\begin{tikzpicture}[scale=0.8]
  \draw[step=1,gray!20,thin] (0,0) grid (6,6);
  \draw[thick,->] (0,0) -- (6.5,0) node[right]{\small $x$};
  \draw[thick,->] (0,0) -- (0,6.5) node[above]{\small $y$};
  \foreach \x in {1,...,6} \draw (\x,0.07) -- (\x,-0.07) node[below]{\tiny \x};
  \foreach \y in {1,...,6} \draw (0.07,\y) -- (-0.07,\y) node[left]{\tiny \y};
  \draw[thick,funBlue] (0,0) -- (6,4);
  \fill[funBlue] (0,0) circle (2.5pt);
  \fill[funBlue] (3,2) circle (2.5pt) node[above left]{\small $(3, 2)$};
  \fill[funBlue] (6,4) circle (2.5pt) node[above left]{\small $(6, 4)$};
\end{tikzpicture}
\end{center}
\end{questionText}
\choiceA{$k = \frac{1}{2}$}
\choiceB{$k = \frac{2}{3}$}
\choiceC{$k = \frac{3}{2}$}
\choiceD{$k = 2$}
\correctAnswer{B}
\explanation{Use any labeled point: $k = \frac{2}{3}$ from $(3, 2)$. Check: $\frac{4}{6} = \frac{2}{3}$ \checkmark.}
\end{practiceQuestion}

\begin{practiceQuestion}{1-3-q27}{gmc}
\begin{questionText}
The table below shows a proportional relationship between gallons of gas and miles driven. What is $k$ and what does it mean?

\begin{center}
\begin{tabular}{|c|c|}
\hline
\textbf{Gallons ($x$)} & \textbf{Miles ($y$)} \\ \hline
$2$ & $54$ \\ \hline
$5$ & $135$ \\ \hline
$8$ & $216$ \\ \hline
\end{tabular}
\end{center}
\end{questionText}
\choiceA{$k = 2$; uses $2$ gallons per trip}
\choiceB{$k = 27$; drives $27$ miles per gallon}
\choiceC{$k = 54$; drives $54$ miles on $2$ gallons}
\choiceD{$k = 135$; drives $135$ miles on $5$ gallons}
\correctAnswer{B}
\explanation{$k = \frac{54}{2} = 27$, $\frac{135}{5} = 27$, $\frac{216}{8} = 27$. The car gets $27$ miles per gallon.}
\end{practiceQuestion}

\begin{practiceQuestion}{1-3-q28}{gmc}
\begin{questionText}
Two proportional lines are graphed below. Line A passes through $(2, 6)$ and Line B passes through $(2, 4)$. Which line has the greater constant of proportionality?

\begin{center}
\begin{tikzpicture}[scale=0.7]
  \draw[step=1,gray!20,thin] (0,0) grid (5,7);
  \draw[thick,->] (0,0) -- (5.5,0) node[right]{\small $x$};
  \draw[thick,->] (0,0) -- (0,7.5) node[above]{\small $y$};
  \foreach \x in {1,...,5} \draw (\x,0.07) -- (\x,-0.07) node[below]{\tiny \x};
  \foreach \y in {1,...,7} \draw (0.07,\y) -- (-0.07,\y) node[left]{\tiny \y};
  \draw[thick,funBlue] (0,0) -- (4,7);
  \fill[funBlue] (2,3.5) circle (2pt);
  \node[above left,funBlue] at (3.5,6.5) {\small Line A};
  \draw[thick,funRed] (0,0) -- (5,7);
  \fill[funRed] (2,2.8) circle (2pt);
  \node[below right,funRed] at (4.5,6.5) {\small Line B};
  \fill[funBlue] (0,0) circle (2pt);
  \fill[funBlue] (2,6) circle (2.5pt) node[above left]{\small $(2,6)$};
  \fill[funRed] (2,4) circle (2.5pt) node[below right]{\small $(2,4)$};
\end{tikzpicture}
\end{center}
\end{questionText}
\choiceA{Line A, with $k = 3$}
\choiceB{Line B, with $k = 2$}
\choiceC{They have the same $k$}
\choiceD{Line B, with $k = 4$}
\correctAnswer{A}
\explanation{Line A: $k = \frac{6}{2} = 3$. Line B: $k = \frac{4}{2} = 2$. Line A has the greater constant of proportionality and is steeper.}
\end{practiceQuestion}

\begin{practiceQuestion}{1-3-q29}{gsa}
\begin{questionText}
The graph shows a proportional relationship. Use the labeled point to find $k$, then determine $y$ when $x = 12$.

\begin{center}
\begin{tikzpicture}[scale=0.6]
  \draw[step=1,gray!20,thin] (0,0) grid (8,8);
  \draw[thick,->] (0,0) -- (8.5,0) node[right]{\small $x$};
  \draw[thick,->] (0,0) -- (0,8.5) node[above]{\small $y$};
  \foreach \x in {2,4,6,8} \draw (\x,0.07) -- (\x,-0.07) node[below]{\tiny \x};
  \foreach \y in {2,4,6,8} \draw (0.07,\y) -- (-0.07,\y) node[left]{\tiny \y};
  \draw[thick,funBlue] (0,0) -- (8,6);
  \fill[funBlue] (0,0) circle (2.5pt);
  \fill[funBlue] (4,3) circle (2.5pt) node[above left]{\small $(4, 3)$};
\end{tikzpicture}
\end{center}
\end{questionText}
\correctAnswer{$k = 0.75$; when $x = 12$, $y = 9$}
\explanation{$k = \frac{3}{4} = 0.75$. Then $y = 0.75 \times 12 = 9$.}
\end{practiceQuestion}

\begin{practiceQuestion}{1-3-q30}{gsa}
\begin{questionText}
The table shows a proportional relationship between minutes and laps completed. Find $k$ and determine how many laps are completed in $15$ minutes.

\begin{center}
\begin{tabular}{|c|c|}
\hline
\textbf{Minutes ($x$)} & \textbf{Laps ($y$)} \\ \hline
$4$ & $3$ \\ \hline
$8$ & $6$ \\ \hline
$12$ & $9$ \\ \hline
$15$ & $?$ \\ \hline
\end{tabular}
\end{center}
\end{questionText}
\correctAnswer{$k = 0.75$; $11.25$ laps in $15$ minutes}
\explanation{$k = \frac{3}{4} = 0.75$. For $15$ minutes: $y = 0.75 \times 15 = 11.25$ laps.}
\end{practiceQuestion}
